\documentclass[11pt,a4paper]{letter}
\usepackage[utf8]{inputenc}
\usepackage[T1]{fontenc}
\usepackage[french]{babel}
\usepackage[left=2.5cm,right=2.5cm,top=2cm,bottom=2cm]{geometry}
\usepackage{hyperref}
\usepackage{xcolor}

% --- Couleurs Safran ---
\definecolor{primary}{RGB}{0, 51, 102}

\hypersetup{colorlinks=true, urlcolor=primary}

\signature{Loïc Aron MBASSI EWOLO}
\address{Loïc Aron MBASSI EWOLO\\
Cergy, France\\
(+33) 7 59 52 33 98\\
\href{mailto:loicaron.mbassiewolo@ensea.fr}{loicaron.mbassiewolo@ensea.fr}\\
\href{https://www.linkedin.com/in/loic-mbassi}{LinkedIn} | \href{https://github.com/Nameless0l}{GitHub}}

\begin{document}

\begin{letter}{Safran Electronics \& Defense\\
Massy}

\opening{Madame, Monsieur,}

Élève-ingénieur en 2A à l'ENSEA (double diplôme avec Polytechnique Yaoundé), je candidate au stage de déploiement multi-plateforme de modèles IA.

Ce stage correspond exactement à ce que je veux faire : j'ai déjà touché aux trois côtés du triangle GPU/NPU/FPGA, et j'ai envie d'aller plus loin.

Côté \textbf{GPU}, j'ai déployé du YOLOv10 sur Nvidia Jetson pour le challenge CoHoMa -- avec les contraintes de latence et de mémoire que ça implique. Je connais PyTorch et l'écosystème ONNX.

Côté \textbf{FPGA}, je travaille actuellement sur un projet SoPC à l'ENSEA avec Quartus et Platform Designer. Je n'ai pas encore utilisé Vitis AI, mais je comprends les contraintes d'implémentation hardware.

Côté \textbf{TinyML/NPU}, j'ai développé un modèle quantifié pour STM32 (projet Harmony Gloves) -- donc je connais les techniques de quantization et les compromis précision/performance.

La partie benchmarking et documentation me parle aussi : au hackathon IndabaX (2e prix), j'ai passé du temps à comparer rigoureusement différentes approches et à documenter les résultats.

Au MILA cet été, j'ai développé des pipelines PyTorch pour analyser le comportement de modèles -- la rigueur scientifique, c'est quelque chose que j'ai appris là-bas.

Le contexte optronique/Défense m'intéresse. Et le fait que ce soit une activité stratégique en lancement, c'est motivant.

Disponible pour un entretien.

\closing{Cordialement,}

\end{letter}
\end{document}
